% Part: turing-machines
% Chapter: machines-computations
% Section: unary-numbers

\documentclass[../../../include/open-logic-section]{subfiles}

\begin{document}

\olfileid[tr]{tur}{mac}{una}
\olsection{Sayıların Tekli Gösterimi}

\begin{explain}
Turing makineleri, şeritlerine yazılmış simge dizileri üzerinde çalışır.
Bir Turing makinesinin kullandığı alfabeye bağlı olarak bu simge dizileri
çeşitli girdi ve çıktıları gösterebilir. Elbette özellikle ilgi çekici
olanlar, \emph{aritmetik} fonksiyonları, yani doğal sayı fonksiyonlarını
hesaplayan Turing makineleridir. Pozitif tam sayıları göstermenin basit bir
yolu, onları tek bir~$\TMstroke$ simgesinin dizileri olarak kodlamaktır.
$n \in \Nat$ ise $n=0$ olduğunda $\TMstroke^n$ boş dizi, aksi durumda ise
tam olarak $n$ tane $\TMstroke$ simgesinden oluşan dizi olsun.
\end{explain}

\begin{defn}[Hesaplama]
Bir Turing makinesi~$M$, $f\colon \Nat^k \to \Nat$ fonksiyonunu ancak ve
ancak $M$
\[
\TMstroke^{n_1} \TMblank \TMstroke^{n_2} \TMblank \dots \TMblank \TMstroke^{n_k}
\]
girdisinde durup $\TMstroke^{f(n_1, \dots, n_k)}$ çıktısını veriyorsa
\emph{hesaplar}.
\end{defn}

\begin{prob}
  Bir Turing makinesi~$M$'nin $f\colon \Nat^k \to \Nat^m$ fonksiyonunu
  hesaplamasının ne anlama geldiğini tanımlayın.
\end{prob}

\begin{ex}\ollabel{ex:adder}
\emph{Toplama:}
$f(n,m)=n+m$ fonksiyonunu hesaplayan bir makine kuralım. Bunun için şeritte
uzunlukları $n$ ve~$m$ olan iki $\TMstroke$ bloğuyla başlayıp $n+m$ tane
$\TMstroke$ simgesinden oluşan tek bir blokla duran bir makine gerekir.
İki girdi $\TMstroke$ bloğu bir~$\TMblank$ ile ayrıldığından bir yöntem,
$\TMblank$ bulunan kareye bir çizgi yazıp son~$\TMstroke$ simgesini
silmektir.
\begin{figure}\[
\begin{tikzpicture}[->,>=stealth',shorten >=1pt,auto,node distance=2.8cm,
                    semithick]
  \tikzstyle{every state}=[fill=none,draw=black,text=black]

  \node[initial,state]         (A)              {$q_0$};
  \node[state]         (B) [right of=A] {$q_1$};
  \node[state]         (C) [right of=B] {$q_2$};

  \path (A) edge node {\TMtrans{\TMblank}{\TMstroke}{\TMstay}} (B)
           edge [loop above] node {\TMtrans{\TMstroke}{\TMstroke}{\TMright}} (B)
        (B) edge [loop above] node {\TMtrans{\TMstroke}{\TMstroke}{\TMright}} (B)
            edge node {\TMtrans{\TMblank}{\TMblank}{\TMleft}} (C)
        (C) edge [loop above] node {\TMtrans{\TMstroke}{\TMblank}{\TMstay}} (C);
\end{tikzpicture}
\]\caption{$f(x,y)=x+y$ fonksiyonunu hesaplayan bir makine}
\ollabel{fig:adder}
\end{figure}
\end{ex}

\begin{prob}
$\tuple{3,2}$ girdisi için \olref[tur][mac][una]{ex:adder} içindeki makinenin
konfigürasyonlarını adım adım izleyin. Makine $0+0$ hesapladığında ne olur?
\end{prob}

\begin{explain}
\olref[rep]{ex:doubler} içinde girdi olarak bir $\TMstroke$ dizisi alan ve
şeritte iki katı sayıda $\TMstroke$ simgesiyle duran bir Turing makinesi,
yani iki katına çıkarıcı makine örneği verdik. Ne var ki çıktı, iki katına
çıkarılmış $\TMstroke$ bloğunun solunda $\TMblank$ simgeleri içerdiğinden
bu makine, varsaymış olabileceğinizin tersine, aslında $f(x)=2x$
fonksiyonunu hesaplamaz. Bunu düzeltmenin iki yolunu açıklayacağız.
\end{explain}

\begin{ex}
\olref{fig:doubler-disc} içindeki makine $f(x)=2x$ fonksiyonunu hesaplar.
\olref[rep]{ex:doubler} içindeki makinenin yaptığı gibi girdideki her
$\TMstroke$ için girdiyi silip en sağa iki $\TMstroke$ yazmak yerine, bu
makine girdideki her~$\TMstroke$ için sağa tek bir~$\TMstroke$ ekler.
Girdinin nerede bittiğini izlemesi gerektiğinden girdiyle eklenen çizgiler
arasında bir~$\TMblank$ bırakır ve en sonunda bu kareyi bir~$\TMstroke$ ile
doldurur. Girdinin neresinde olduğumuzu da ``anımsamamız'' gerekir; bu yüzden
girdi bloğundaki bir~$\TMstroke$ simgesini geçici olarak bir~$\TMblank$ ile
değiştiririz.
  \begin{figure}
\[
\begin{tikzpicture}[->,>=stealth',shorten >=1pt,auto,node distance=2.8cm,
                    semithick]
  \tikzstyle{every state}=[fill=none,draw=black,text=black]

  \node[initial,state] (0)              {$q_0$};
  \node[state]         (1) [above of=0] {$q_1$};
  \node[state]         (2) [above of=1] {$q_2$};
  \node[state]         (3) [right of=2] {$q_3$};
  \node[state]         (4) [below of=3] {$q_4$};
  \node[state]         (5) [below of=4] {$q_5$};
  \node[state]         (6) [above right of=4] {$q_6$};
  \node[state]         (7) [below of=6] {$q_7$};
  \node[state]         (8) [below of=7] {$q_8$};

  \path (0) edge node {\TMtrans{\TMstroke}{\TMblank}{\TMright}} (1)
%    (0) edge node {\TMtrans{\TMblank}{\TMblank}{\TMstay}} (h)
    (1) edge [loop left] node {\TMtrans{\TMstroke}{\TMstroke}{\TMright}} (1)
      edge node {\TMtrans{\TMblank}{\TMblank}{\TMright}} (2)
    (2) edge [loop above] node {\TMtrans{\TMstroke}{\TMstroke}{\TMright}} (2)
        edge node {\TMtrans{\TMblank}{\TMstroke}{\TMleft}} (3)
    (3) edge [loop above] node {\TMtrans{\TMstroke}{\TMstroke}{\TMleft}} (3)
        edge node[left] {\TMtrans{\TMblank}{\TMblank}{\TMleft}} (4)
    (4) edge        node[left] {\TMtrans{\TMstroke}{\TMstroke}{\TMleft}} (5)
    (5) edge [loop below] node {\TMtrans{\TMstroke}{\TMstroke}{\TMleft}} (5)
        edge              node {\TMtrans{\TMblank}{\TMstroke}{\TMright}} (0)
    (4) edge      node[sloped] {\TMtrans{\TMblank}{\TMstroke}{\TMright}} (6)
    (6) edge              node {\TMtrans{\TMblank}{\TMstroke}{\TMright}} (7)
    (7) edge [loop right] node {\TMtrans{\TMstroke}{\TMstroke}{\TMright}} (7)
        edge              node {\TMtrans{\TMblank}{\TMblank}{\TMleft}} (8)
    (8) edge [loop right] node {\TMtrans{\TMstroke}{\TMblank}{\TMstay}} (8);
    \end{tikzpicture}
\]
\caption{$f(x)=2x$ fonksiyonunu hesaplayan bir makine}
\ollabel{fig:doubler-disc}
\end{figure}
\end{ex}

\begin{ex}\ollabel{ex:mover}
$f(x)=2x$ hesaplamanın ikinci bir yolu, özgün iki katına çıkarıcı makineyi
koruyup sonuna iki katına çıkarılmış çizgi bloğunu şeridin en soluna taşıyan
durumlar ve yönergeler eklemektir. \olref{fig:mover} içindeki makine yalnızca
bu son kısmı yapar: Bir $\TMblank$ bloğunu izleyen bir $\TMstroke$ bloğundan
oluşan şeritte başlatıldığında (kafa $\TMblank$ bloğunun herhangi bir yerinde
olabilir), $\TMstroke$ simgelerini birer birer silip şeridin başına yazar.
Ne zaman bittiğini anlayabilmek için önce $\TMstroke$ bloğunun sonunu bir
$\TMendtape$ simgesiyle işaretler; bu simge sonunda silinir. Durumları
$q_6$'dan başlayarak numaralandırdık, böylece iki katına çıkarıcı makineye
eklenebilirler. Gerekli olan tek şey ek bir $\delta(q_0,
\TMblank)=\tuple{q_6,\TMblank,\TMstay}$ yönergesi, yani $q_0$'dan $q_6$'ya
$\TMtrans{\TMblank}{\TMblank}{\TMstay}$ etiketli bir oktur. (İnce bir sorun
vardır: Ortaya çıkan makine $x=0$ girdisinde çalışmaz. Bunu alıştırma olarak
bırakıyoruz.)
\begin{figure}
  \[
  \begin{tikzpicture}[->,>=stealth',shorten >=1pt,auto,node distance=2.8cm,
                      semithick]
    \tikzstyle{every state}=[fill=none,draw=black,text=black]
    \node[initial,state] (6)              {$q_6$};
    \node[state]         (7) [right of=6] {$q_7$};
    \node[state]         (8) [right of=7] {$q_8$};
    \node[state]         (9) [below of=8] {$q_9$};
    \node[state]         (10) [left of=9] {$q_{10}$};
    \node[state]         (11) [left of=10]  {$q_{11}$};
    \node[state]         (12) [below of=11]  {$q_{12}$};
    \node[state]         (13) [right of=12] {$q_{13}$};
    \node[state]         (14) [right of=13] {$q_{14}$};
    \path
    (6)  edge node {\TMtrans{\TMblank}{\TMblank}{\TMright}} (7)
    (7)  edge [loop above] node {\TMtrans{\TMblank}{\TMblank}{\TMright}} (7)
         edge node {\TMtrans{\TMstroke}{\TMstroke}{\TMright}} (8)
    (8)  edge [loop above] node {\TMtrans{\TMstroke}{\TMstroke}{\TMright}} (8)
         edge node {\TMtrans{\TMblank}{\TMendtape}{\TMleft}} (9)
    (9)  edge [loop right] node {\TMtrans{\TMstroke}{\TMstroke}{\TMleft}} (9)
         edge node {\TMtrans{\TMblank}{\TMblank}{\TMright}} (10)
    (10) edge node[above] {\TMtrans{\TMstroke}{\TMblank}{\TMleft}} (11)
    (11) edge [loop above] node {\TMtrans{\TMblank}{\TMblank}{\TMleft}} (11)
         edge node[left] {\begin{tabular}{@{}l@{}}
          \TMtrans{\TMendtape}{\TMendtape}{\TMright}\\
          \TMtrans{\TMstroke}{\TMstroke}{\TMright}
         \end{tabular}} (12)
    (12) edge node[] {\TMtrans{\TMblank}{\TMstroke}{\TMright}} (13)
    (13) edge [loop below] node {\TMtrans{\TMblank}{\TMblank}{\TMright}} (13)
         edge node[sloped] {\TMtrans{\TMstroke}{\TMblank}{\TMleft}} (11)
    (13) edge node {\TMtrans{\TMendtape}{\TMblank}{\TMstay}} (14);
  \end{tikzpicture}
  \]
  \caption{Bir $\TMstroke$ bloğunu sola taşıma}
  \ollabel{fig:mover}
\end{figure}
\end{ex}

\begin{prob}
\olref[tur][mac][una]{ex:mover} içinde
\olref[tur][mac][una]{fig:doubler-disc} içindeki iki katına çıkarıcı makineyle
\olref[tur][mac][una]{fig:mover} içindeki taşıyıcı makinenin birleşiminden
oluşan bir makine betimledik. Bu birleşik makineyi $x=0$ girdisinde, yani boş
bir şeritte başlatırsanız ne olur? Makineyi, bu durumda $2x=0$ çıktısıyla
duracak biçimde nasıl düzeltirdiniz? (Bunu bir durum ve bir geçiş ekleyerek
yapabilmelisiniz.)
\end{prob}

\begin{prob}
\emph{Çıkarma:} Uzunlukları $n$ ve~$m$ olan, $n>m$ koşulunu sağlayan ve boş
olmayan iki çizgi dizgesi girdi olarak verildiğinde $f(n,m)=n-m$ fonksiyonunu
hesaplayan bir Turing makinesi tasarlayın.
\end{prob}

\begin{prob}
\emph{Eşitlik:} Aşağıdaki fonksiyonu hesaplayan bir Turing makinesi
tasarlayın:
\[
\fn{equality}(n,m) = 
\begin{cases}
  \text{1} & \text{$n = m$ ise} \\
  \text{0} & \text{$n \neq m$ ise}
\end{cases}
\]
burada $n$ ve~$m \in \PosInt$.
\end{prob}

\begin{prob}
$x$ ve $y$, şeritte aralarında bir $\TMblank$ bulunan $\TMstroke$
dizgeleriyle gösterilen pozitif tam sayılar olmak üzere $\min(x,y)$
fonksiyonunu hesaplayan bir Turing makinesi tasarlayın. Makinenin
alfabesinde ek simgeler kullanabilirsiniz.

$\min$ fonksiyonu argümanlarının en küçük değerini seçer; dolayısıyla
$\min(3,5)=3$, $\min(20,16)=16$, $\min(4,4)=4$ vb.
\end{prob}

\begin{defn}
  Bir Turing makinesi~$M$, $f\colon \Nat^k \pto \Nat$ kısmi fonksiyonunu
  ancak ve ancak
  \begin{enumerate}
    \item $f(n_1, \dots, n_k)=m$ ise $M$, $\TMstroke^{n_1}\concat\TMblank\concat
    \dots \concat\TMblank\concat\TMstroke^{n_k}$ girdisinde
    $\TMstroke^{m}$ çıktısıyla duruyorsa;
    \item $f(n_1, \dots, n_k)$ tanımsızsa $M$ hiç durmuyor ya da tek bir
    $\TMstroke$ bloğu olmayan bir çıktıyla duruyorsa
  \end{enumerate}
  \emph{hesaplar}.
\end{defn}

\end{document}
